\section{Recomendaciones}

    \begin{itemize}
        \item Validar el aplicativo en un entorno industrial real, involucrando personal de mantenimiento durante la ejecución de actividades preventivas. Esto permitiría evaluar variables como la usabilidad, la aceptación por parte de los usuarios, la reducción de tiempos de intervención y el impacto en la calidad de los procedimientos.
        \item Comparar el enfoque basado en marcadores QR con otros métodos de activación de realidad aumentada, como el reconocimiento de objetos, superficies o el posicionamiento manual de modelos, con el fin de determinar cuál ofrece mejores resultados según el tipo de maquinaria y las condiciones del entorno.
        \item Ampliar la evaluación incorporando diferentes tipos de dispositivos móviles y escenarios de iluminación, distancia y oclusión, para analizar el comportamiento del sistema bajo condiciones más cercanas a las encontradas en aplicaciones industriales.
        \item Explorar la incorporación de indicadores propios de la ingeniería de mantenimiento, tales como MTTR, MTBF, disponibilidad operacional o modelos de predicción de fallas, con el propósito de complementar la asistencia visual con herramientas de apoyo a la toma de decisiones.
        \item Incorporar un sistema híbrido de almacenamiento que permita sincronizar la información con un servidor cuando exista conectividad y mantener el funcionamiento sin conexión mediante almacenamiento local, garantizando la disponibilidad de la información en entornos industriales.
        \item Implementar un sistema de usuarios y roles que diferencie las funciones de administradores, supervisores y ejecutores de mantenimiento, restringiendo la modificación de información técnica únicamente al personal autorizado.
        \item Agregar comandos de voz para consultar o editar la información técnica de los equipos de mantenimiento, cuando el operario utilice elementos de protección personal que le dificulten la manipulación del dispositivo.
        \item Ampliar el soporte para diferentes tipos de marcadores o técnicas de seguimiento distintas a los códigos QR, permitiendo adaptar el sistema a diversos escenarios industriales.
        \item Hacer uso de las características de realidad extendida de SolidWorks para la exportación de los modelos en el estándar glTF reduciendo el tiempo requerido para preparar nuevos modelos tridimensionales.
    \end{itemize}